% !TEX program = lualatex
\documentclass[a4paper,11pt]{scrartcl}

\usepackage[
  lang=en,
  mode=journal,
  theme=graphite,
  toc=true,
  toc-layout=page
]{synaptic}
\usepackage{booktabs,tabularx}
\newcolumntype{Y}{>{\raggedright\arraybackslash}X}
\newcommand{\DocCmd}[1]{\textcolor{syn-primary!80!black}{\texttt{\textbackslash#1}}}

\SynapticHeader{Synaptic Technical Reference}
\SynapticSubHeader{Architecture, invariants, and release workflow}
\SynapticPreprint{v5.0.0}
\SynapticTitle{synaptic Technical Reference}
\SynapticAuthor{synaptic project}
\SynapticAffiliation{Open-source LaTeX community}
\SynapticSubject{Technical reference for synaptic version 5.0.0}
\SynapticKeywords{LaTeX3, architecture, l3build, regression testing}

\begin{synabstract}
  This reference describes synaptic's module graph, load-time invariants,
  namespace rules, font and colour subsystems, extension points, generated-file
  policy, automated verification workflow, and packaging details.
\end{synabstract}

\begin{document}
\SynapticMakeTitle

\section{Architecture}

\begin{verbatim}
synaptic.sty
  `-- synaptic-base.sty
      |-- synaptic-lang-{en,zh}.def
      |-- synaptic-color.sty
      |-- synaptic-fonts.sty
      |-- synaptic-layout.sty
      |-- synaptic-theorem.sty
      |-- synaptic-title.sty       shared metadata/renderers
      `-- mode module(s)
          |-- synaptic-journal.sty     journal
          |-- synaptic-book.sty        book
          |-- synaptic-lecture.sty     lecture
          `-- synaptic-notes.sty       notes
\end{verbatim}

The thin shell forwards user options to \texttt{synaptic-base} and provides the
namespaced acknowledgments command. The base module checks the engine and
class, parses options, loads labels, validates structural requirements, and
dispatches core and mode-specific modules. Language definitions load before
the modules that consume them.

\subsection{Module responsibilities}

\noindent{\setlength{\parindent}{0pt}\noindent\begin{tabularx}{\textwidth}{@{}l l Y@{}}
  \toprule
  \textbf{Module} & \textbf{Kind} & \textbf{Responsibility} \\
  \midrule
  \texttt{synaptic.sty} & shell & option forwarding, acknowledgments \\
  \texttt{synaptic-base.sty} & core & engine/class checks, keys, dispatch, setup \\
  \texttt{synaptic-lang-*.def} & data & bilingual label constants \\
  \texttt{synaptic-color.sty} & core & theme registry, semantic tokens, delivery \\
  \texttt{synaptic-fonts.sty} & core & fixed Latin and CJK font stack \\
  \texttt{synaptic-layout.sty} & core & geometry, typography, headings, numbering \\
  \texttt{synaptic-theorem.sty} & core & theorem styles, environments, proofs \\
  \texttt{synaptic-title.sty} & core & metadata model and four title renderers \\
  \texttt{synaptic-journal.sty} & mode & journal navigation and statements \\
  \texttt{synaptic-book.sty} & mode & chapter/part style, matter helpers \\
  \texttt{synaptic-lecture.sty} & mode & masthead, teaching cards, navigation \\
  \texttt{synaptic-notes.sty} & mode & notes masthead and knowledge cards \\
  \bottomrule
\end{tabularx}}

\noindent{\setlength{\parindent}{0pt}\noindent\begin{tabularx}{\textwidth}{@{}l l Y@{}}
  \toprule
  \textbf{Mode} & \textbf{Class invariant} & \textbf{Extra modules} \\
  \midrule
  journal & KOMA article/report & journal \\
  book & chapter-based KOMA class with front matter & book \\
  lecture & KOMA article/report & lecture \\
  notes & KOMA article/report & notes \\
  \bottomrule
\end{tabularx}}

\section{Configuration lifecycle}

The package-level \texttt{synaptic} key family is evaluated before module
loading. Mode, language, title layout, measure, density, numbering,
delivery colour mode, theorem style, and binding offset therefore become
structural invariants. The runtime key family
exposes only the state that can be applied without reloading packages:
\texttt{theme}, \texttt{toc}, and \texttt{toc-layout}. Structural or unknown
keys in a runtime request are errors; unknown package options are errors too.

Book mode is rejected unless both chapter and front-matter interfaces exist.
LuaLaTeX and KOMA-Script requirements are fatal diagnostics rather than delayed
undefined-control-sequence failures.

\section{Design invariants}

The following properties are deliberately asserted by the implementation and
protected by the regression suite.

\begin{itemize}
  \item \textbf{Deterministic typography.} A fixed font stack and explicit
    spacing, measure, and heading sizes prevent machine-dependent redesigns.
  \item \textbf{No global bad-box suppression.} Emergencystretch and tolerance
    are tuned; overfull and underfull warnings remain visible for the author.
  \item \textbf{No forced paper size or page-number reset.} Geometry follows
    the class; journal titles do not silently restart numbering.
  \item \textbf{One Unicode conversion.} PDF metadata is passed to hyperref as
    source tokens; pre-conversion would double-encode.
  \item \textbf{Owned public names.} Every public command and environment is
    namespaced; generic environments supplied by other packages stay untouched.
  \item \textbf{Semantic colour architecture.} Components consume semantic
    tokens; themes and delivery modes change tokens, never component logic.
\end{itemize}

\section{Subsystems}

\subsection{Colour tokens}

Built-in themes are stored in a global property list. Activating a theme maps
its primary, secondary, and accent colours to semantic public names and derives
surface, border, rule, link, readable-muted, decorative-muted, and text tokens.
Print and mono delivery modes alter derived roles without changing component
code. Danger, warning, and success
colours remain semantic constants so their meaning does not change with brand
colour.

\begin{verbatim}
\SynapticDefineTheme{name}{primary}{secondary}[accent]
\SynapticSetup{theme=name}
\end{verbatim}

The colour arguments are xcolor expressions. Theme identifiers are validated
lowercase slugs, definitions are immutable, and activation is explicit.

\subsection{Font selection}

Font discovery only validates the documented fixed stack. Latin text and
mathematics use Libertinus Serif, Sans, Mono, and Math; the bold mathematics
version is a deterministic synthetic weight of the same Libertinus Math face.
For \texttt{lang=zh},
ctex is loaded without a preset font scheme; Source Han Serif/Sans CN provide
Chinese text and LXGW WenKai provides Chinese emphasis. Missing fonts are
package errors rather than fallback triggers.

\subsection{Layout and document structure}

Geometry combines mode defaults with explicit measure and density profiles.
Journal and lecture favour moderate line lengths; notes supports compact
single/two-column composition; book uses mirrored inner/outer margins, binding
correction, outer running heads, and empty duplex blanks. Widow and club penalties,
vertical rhythm, list spacing, line spacing, and paragraph indentation are
explicit rather than hidden behind global bad-box suppression.

KOMA heading fonts and skips are installed in the
\texttt{begindocument/before} hook, with all numbered levels sharing the sans
heading family: the primary colour for \texttt{section}, dark for
\texttt{subsection}, and a smaller, muted tint for \texttt{subsubsection} so
the depth hierarchy stays legible. Hyperref, bookmark, microtype, captions,
lists, caption punctuation, and mode-aware equation numbering are configured centrally.
A hairline running-head separator (\texttt{headsepline}) is installed by the
mode modules through a shared helper and coloured by the dedicated KOMA
\texttt{headsepline} font, keeping decoration lighter than head text.

The table of contents is refined through KOMA's \texttt{\string\DeclareTOCStyleEntry}
immediately before it is typeset: top-level entries repeat the on-page heading
colour in bold, deeper entries remain dark, and number alignment and vertical
rhythm are tightened. Figure and table caption labels carry the secondary
theme token.

\subsection{Titles and metadata}

Title values are held in global expl3 token lists and sequences because they
are document-wide state. Authors and their
affiliation marks use parallel sequences; affiliations and emails are
repeatable. A common metadata model feeds four renderers: inline journal,
book title sequence, course masthead, and compact notes masthead. The title
command validates a non-empty title and controls optional TOC placement without
silently resetting the user's page numbering. The namespaced
\texttt{synabstract} environment captures abstract text; standard
\texttt{abstract} and \verb|\maketitle| remain untouched.

PDF information is passed as source tokens to hyperref. Hyperref performs the
Unicode PDF-string conversion once; pre-converting and then calling
\verb|\hypersetup| would double-encode metadata.

\subsection{Statements}

thmtools defines plain, definition, and remark styles on top of amsthm. The
built-in statement family (\texttt{syntheorem}, \texttt{synlemma},
\texttt{synproposition}, \texttt{syndefinition}, \texttt{syncorollary}, \texttt{synaxiom},
\texttt{synconjecture}, \texttt{synclaim}, \texttt{synproperty}, \texttt{synexample},
\texttt{synexercise}, \texttt{syncriterion}, \texttt{synconvention},
\texttt{synproblem}, \texttt{synsolution}, \texttt{synremark}, \texttt{synnote}, and
\texttt{synwarning}) uses a mode-aware shared counter: chapter scope for books and
section scope elsewhere by default. Under \texttt{theorem-style=boxed} the
numbered statement environments are wrapped with tcolorbox (a theme-tinted
background with a west accent bar) at \texttt{\string\begin{document}};
\texttt{synproof} stays unboxed and the standard \texttt{proof} is not
redefined. \texttt{theorem-style=plain} disables statement wrapping. Examples
and exercises are plain-style statements; warnings use the remark presentation.
Because every statement shares the theorem counter, the module rebuilds each
environment's hyperlink anchor (\texttt{theH<name>}) against the real structural
parent so links stay unique across sections and chapters.

\section{Namespace policy}

\noindent{\setlength{\parindent}{0pt}\noindent\begin{tabularx}{\textwidth}{@{}l l Y@{}}
  \toprule
  \textbf{Scope} & \textbf{Pattern} & \textbf{Example} \\
  \midrule
  Public command & \texttt{\textbackslash Synaptic\ldots} & \DocCmd{SynapticDefineTheme} \\
  Public environment & \texttt{syn...} & \texttt{synsummary} \\
  Public colour token & \texttt{syn-...} & \texttt{syn-primary} \\
  Internal function & \texttt{\textbackslash\_\_synaptic\_\ldots:} & title rendering helper \\
  Private local state & \texttt{\textbackslash l\_\_synaptic\_\ldots} & title-layout scratch value \\
  Private global state & \texttt{\textbackslash g\_\_synaptic\_\ldots} & mode, metadata, or registry \\
  Private constant & \texttt{\textbackslash c\_\_synaptic\_\ldots} & label or environment family \\
  \bottomrule
\end{tabularx}}

The expl3 prefixes \texttt{g} and \texttt{l} describe lifetime, not ownership.
Ownership is carried independently by the \texttt{synaptic} module name and
the private double underscore. Document-wide state is therefore still private
\texttt{g\_\_synaptic\_...} state; short-lived scratch values use
\texttt{l\_\_synaptic\_...}.

New public names must follow the same ownership rules. Custom statement
environment names accepted by \verb|\SynapticNewTheorem| begin with
\texttt{syn} and contain only lowercase letters and digits; duplicate or
unnamespaced names are errors.

\section{Extension points}

\begin{itemize}
  \item Add labels in a new \texttt{synaptic-lang-*.def} file, then extend the
    language key and definition-file dispatch.
  \item Treat a font-stack change as an intentional design and dependency
    migration; update validation, both manuals, and representative documents.
  \item Add modes in the base key family, layout dispatch, and mode-module
    loader. State the required KOMA class capability and test it.
  \item Build visual extensions from semantic colour tokens rather than fixed
    RGB values.
  \item Add tests along with every extension; the suite is the guard for the
    invariants above.
\end{itemize}

\section{Automated verification}

\texttt{synaptic.dtx} is the source of truth. DocStrip uses
\texttt{synaptic.ins} to generate the installable \texttt{.sty} and
\texttt{.def} files under \texttt{build/unpacked}. Extracted modules, compiled
PDFs, and release archives are deliberately untracked, so there is only one
implementation source to review and update.

The regression suite is organised as follows:

{\setlength{\parindent}{0pt}\noindent\begin{tabularx}{\textwidth}{@{}l Y@{}}
  \toprule
  \textbf{Group} & \textbf{Coverage} \\
  \midrule
  \texttt{journal}, \texttt{lecture}, \texttt{notes}, \texttt{book} & core
  mode behaviour, navigation, and defaults \\
  \texttt{book-zh}, \texttt{lecture-zh}, \texttt{journal-zh},
  \texttt{language-zh} & Chinese labels, CJK fonts, and bilingual layering \\
  \texttt{themes} & cycling all five themes and a custom theme \\
  \texttt{color-mode-print}, \texttt{options} & print and monochrome delivery \\
  \texttt{customization} & runtime setup, custom theorems \\
  \texttt{toc-layouts}, \texttt{title-layouts} & TOC and title compositions \\
  \texttt{layout-profiles} & density, measure, binding offset \\
  \texttt{theorem-boxes}, \texttt{strict-options} & full statement family,
  namespace validation, and strict option errors \\
  \texttt{book-full} & complete title sequence and matter workflow \\
  \texttt{lecture-continuous} & continuous numbering profile \\
  \bottomrule
\end{tabularx}}

Each group typesets a document with the standard engine and compares the
normalised log against its committed expectation file. The shared
\texttt{scripts/verify} entry point runs the suite, builds the documented
source, four manuals, and all examples, then fails on warnings, overfull boxes,
missing characters, or TeX errors. CI invokes the same entry point.

\section{Source and release workflow}

\begin{verbatim}
./scripts/verify  # tests, manuals, examples, diagnostics
l3build unpack    # extract .sty/.def into build/unpacked
l3build ctan      # create tested TDS and CTAN archives
l3build install   # copy package files into the texmf tree
l3build uninstall # remove a previous installation
\end{verbatim}

The \texttt{ctan} target emits \texttt{synaptic-ctan.zip} (the source tree
plus generated documentation and examples) together with the nested
\texttt{synaptic.tds.zip} in TDS layout. Installation replaces the previous
version in place. The repository does not retain compatibility copies or
migration shims; the changelog records removed interfaces.

\end{document}
